2.- Utilizaremos dole contención: \\

\underline{$\subseteq$}

P.D para toda cadena $w\in L$, se cumple $w=w^{R}$

Utilizando Inducción:

-Caso Base:
\begin{itemize}


\item $w =\lambda$:

Por definición recursiva de la reversa, sabemos que $\lambda^{R}=\lambda$

Por lo tanto se cumple

$\lambda = \lambda^{R}$

\item $(w = a \mid a \in\Sigma)$:

por definicón recursiva de la reversa:

$$a^{R}=(\lambda a)^{R}=a\cdot\lambda^{R}=a\cdot\lambda=a$$

por lo tanto $a = a^{R}$
\end{itemize}
-Hipótesis de Inducción:

Supongamos que para toda cadena $w\in L$, se cumple:

$$w=w^{R}$$

-Paso Inductivo:

P.D que: 
$$(awa)^{R}=awa$$

por la propiedad de la reversa de una concatenación $(\alpha\beta\gamma)^{R}=\gamma^{R}\beta^{R}\alpha^{R}$ 

obtenemos:

$$((awa)^{R}=(a^{R})\cdot(w^{R})\cdot(a^{R}))$$

Aplicando el caso base $(a^{R}=a)$:

$$a^{R}\cdot w^{R}\cdot a^{R}=a\cdot w^{R}\cdot a$$

Aplicando la hipótesis de inducción $(w^{R}=w)$:
$$a\cdot w\cdot a=awa$$

Por lo tanto 
$$(awa)^{R}=awa$$

Se cumple que 

$$L\subseteq\{w\in\Sigma^{*}\mid w=w^{R}\}$$

\medskip
\underline{$\supseteq$}

P.D que $\{w\in\Sigma^{*}\mid w=w^{R}\}\subseteq L$

Si $w\in\Sigma^{*}$ es palindromo (es decir $w=w^{R}$), entonces $w\in L$.

Haremos inducción sobre la longitud de $w$ ($|w|= n$), ya que el conjunto de palíndromos está definida por propiedades (como la longitud de la cadena) y no por una definición recursiva.

-Caso Base:
\begin{itemize}
\item $|w|=0$: \\
Si $|w|=0$ la unica cadena que cuumple con esta propiedad es 
$$w=\lambda$$

\item $|w|=\lambda$: \\
Por definción de $L$, sabemos que $\lambda\in L$

\item $|w|=1$: \\
Si $|w|=1$, entonces 
$$w=a$$
para $a\in\Sigma$.

Por definción de $L$, sabemos que $a\in L$ 

\end{itemize}

-Hipótesis de Inducción:

Supongamos que para roda cadena $w\in\Sigma^{*}$ con longitud:
$ 0 \leq |w| \leq n$, con $n \geq 2$, si w es un palíndromo $(w =w^{R})$, entonces $w\in L$.

-Paso Inductivo:

Sea $w\in\Sigma^{*}$ con $|w|= n\geq 2$ tal que $w=w^{R}$, w es un palíndromo

P.D que $w\in L$

Como $|w| \geq 2$, podemos escribir:
$$w = \alpha \cdot u \cdot \beta$$ donde $\alpha,\beta\in\Sigma$ y $u\in\Sigma^{*}$

Además $|w| = |\alpha| + |u| + |\beta|$, es decir:
$$|w| = 1 + |u| + 1$$
$$|w| = |u| + 2$$
$$|u| = |w| - 2$$
$$|u| = n - 2$$

Aplicando $(w = w^{R})$ obtenemos:
$$\alpha \cdot u \cdot \beta = (\alpha \cdot u \cdot \beta)^{R}= \beta^{R} \cdot u^{R} \cdot \alpha^{R}$$
$$= \beta \cdot u^{R} \cdot \alpha$$

Igualando los extremos de la cadena obtenemos:
$$\alpha = \beta$$
$$u = u^{R}$$

Podemos deducir que:
\begin{itemize}
    \item $w$ es de la forma $w = \alpha \cdot u \cdot \alpha$ 
    \item $u$ es un palíndromo, ya que $u = u^{R}$ y $|u| = n - 2 < n$
\end{itemize}

Por hipótesis de inducción, $u\in L$ 
Como $u \in L$ y $\alpha \in \Sigma$, por el tercer punto de la definción de $L$:
$$\alpha \cdot u \cdot \alpha \in L$$
Es decir:
$$w \in L$$

Al haber demostrado ambas contenciones, concluimos que:
$$L = \{w\in\Sigma^{*}\mid w=w^{R}\}$$